\section{Proved limits}\label{sec:negative}

Three statements of the preceding sections are limits, and each is proved.
\begin{itemize}
\item \emph{The dispersion series does not converge beyond $k_*$} (Theorem~\ref{thm:radius}(c)). A series with a larger radius would make $\alpha_h$ analytic at $x_*$, where it has a square-root branch point. The negative control confirms this: the same tiling, asked to cover the disk of radius $0.158>x_*$, fails on the real axis beyond $x_*$, where the two zeros are a conjugate pair.
\item \emph{Beyond the collision the shear modes propagate} (Theorem~\ref{thm:radius}(d)). For real $x$ slightly above $x_*$ the two roots near $\alpha_*$ are complex conjugates, so the purely damped regime ends exactly at $k_*$.
\item \emph{The two-pole model misplaces the collision.} The telegraph model $\alpha^2+\alpha+2x=0$, calibrated to the diffusion constant and to the first non-hydrodynamic mode, collides at $(q^2,\alpha)=(\frac12,-\frac12)$, while the exact collision is at $(0.6060,-0.5697)$. So the radius is a property of the exact dispersion relation and is not captured by the two-pole truncation.
\end{itemize}
The rates of Section~\ref{sec:rates} carry a limit of scope that the workbench states itself: the zero-quotient sequence of the map to Yang--Mills exists only along couplings $g_j\to0$, and it rests on (H1)--(H4).
